% Generated draft: requires mathematical and source review.
\documentclass[11pt]{article}
\usepackage[margin=25mm]{geometry}
\usepackage{fontspec}
\setmainfont{Latin Modern Roman}
\usepackage{amsmath,amssymb,mathtools,unicode-math}
\setmathfont{Latin Modern Math}
\usepackage{enumitem,needspace,xurl,hyperref}
\setlength{\parindent}{0pt}
\setlength{\parskip}{5pt}
\setlength{\emergencystretch}{3em}
\newcommand{\N}{\mathbb{N}}
\newcommand{\Z}{\mathbb{Z}}
\newcommand{\Q}{\mathbb{Q}}
\newcommand{\R}{\mathbb{R}}
\newcommand{\C}{\mathbb{C}}
\newcommand{\F}{\mathbb{F}}
\newcommand{\A}{\mathbb{A}}
\newcommand{\PP}{\mathbb{P}}
\newcommand{\E}{\mathbb{E}}
\newcommand{\eps}{\varepsilon}
\newcommand{\dd}{\,\mathrm{d}}
\begin{document}
\section*{UnsolvedMath ID: 30001478}
\textbf{Source problem code:} catalogue entry
\subsection{Definitions and mathematical statement}
Let \(k\) be an uncountable algebraically closed field, let \(k_0\) be its prime subfield, and let \(\alpha,\beta\in k\) be algebraically independent over \(k_0\); that is, \(P(\alpha,\beta)\ne0\) for every nonzero \(P\in k_0[u,v]\). No characteristic restriction is supplied. Form the associative unital algebra
\[
R=R(\alpha,\beta)=k\langle x_1,x_2,x_3,x_4\rangle/(f_1,\ldots,f_6),
\]
where the denominator denotes the two-sided ideal generated by
\[
\begin{aligned}
f_1&=x_1(\alpha x_1-x_3)+x_3(x_1-\alpha x_3),\\
f_2&=x_1(\alpha x_2-x_4)+x_3(x_2-\alpha x_4),\\
f_3&=x_2(\alpha x_1-x_3)+x_4(x_1-\alpha x_3),\\
f_4&=x_2(\alpha x_2-x_4)+x_4(x_2-\alpha x_4),\\
f_5&=x_1(\beta x_1-x_2)+x_4(x_1-\beta x_2),\\
f_6&=x_1(\beta x_3-x_4)+x_4(x_3-\beta x_4).
\end{aligned}
\]
Give each generator degree one, so that \(R=\bigoplus_{n\ge0}R_n\) and \(R_0=k\). The augmentation homomorphism \(\varepsilon:R\to k\) sends every \(x_i\) to zero, and its kernel is
\[
R_+=\bigoplus_{n\ge1}R_n.
\]
The original supplied statement describes this algebra as birationally commutative and of Gelfand--Kirillov dimension \(4\); these are retained source descriptions, not independently verified conclusions here.

A two-sided ideal \(P\) is graded if \(P=\bigoplus_{n\ge0}(P\cap R_n)\). It is prime if \(P\ne R\) and, for all two-sided ideals \(I,J\subseteq R\), the inclusion \(IJ\subseteq P\) implies \(I\subseteq P\) or \(J\subseteq P\).

For every field and parameter pair satisfying the stated hypotheses, the first question asks whether \(R_+\) is the only nontrivial graded prime ideal. Under the conventional interpretation that ``nontrivial'' means nonzero and proper, its precise formulation is
\[
\{P\subsetneq R: P\ne(0),\ P\text{ is a graded two-sided prime ideal}\}=\{R_+\}.
\]
The source does not explicitly define ``nontrivial''; the displayed formulation records this interpretation rather than resolving that omission.

The second question asks whether \(R\) admits no nontrivial map to any twisted homogeneous coordinate ring. To explain the terminology, a customary construction starts with a projective \(k\)-scheme \(X\), an automorphism \(\sigma:X\to X\), and an invertible sheaf \(\mathcal L\). Set
\[
\mathcal L_0=\mathcal O_X,\qquad
\mathcal L_n=\mathcal L\otimes\sigma^*\mathcal L\otimes\cdots\otimes(\sigma^{n-1})^*\mathcal L\quad(n\ge1),
\qquad
B(X,\mathcal L,\sigma)=\bigoplus_{n\ge0}H^0(X,\mathcal L_n).
\]
For homogeneous sections \(s\) of degree \(m\) and \(t\) of degree \(n\), multiplication is induced by \(s\otimes(\sigma^m)^*t\). This construction explains a conventional meaning of the target terminology; the supplied record does not specify the intended class of \(X,\mathcal L,\sigma\), including any integrality, dimension, or ampleness requirements.

The record also does not specify whether maps must be unital, \(k\)-linear, or grading preserving, or what makes a map ``nontrivial.'' For example, for a unital \(k\)-algebra target \(B\), the augmentation gives a map \(R\xrightarrow{\varepsilon}k\to B\); the record does not explicitly say whether such maps are excluded. Consequently the second question is retained as a source conjectural target, but a fully quantified assertion about maps cannot be fixed without additional source clarification.
\subsection{Short English statement}
For every stated field and algebraically independent parameter pair, the problem asks whether the augmentation ideal is the unique nonzero proper graded prime ideal and whether there are no nontrivial maps to twisted homogeneous coordinate rings. The meaning of nontrivial maps and the permitted target rings are unspecified. Both questions are supplied as conjectures requiring human mathematical review.
\subsection{Sources}
\begin{itemize}
\item Local UnsolvedMath registry: supplied complete record for permanent ID \(30001478\), exact source problem code \texttt{catalogue entry}, titled ``Prime Ideals and Coordinate-Ring Maps of \(R(\alpha,\beta)\),'' in the Oberwolfach Reports Open Problems collection. The supplied proposed year is 2010; no individual proposer is identified.
\item Supplied primary source citation: \emph{Interactions between Algebraic Geometry and Noncommutative Algebra} (2010), Oberwolfach Reports, pp.\ 1319--1381, DOI \texttt{10.4171/owr/2010/22}. \url{https://doi.org/10.4171/owr/2010/22}. The collection describes these reports as documenting mathematical workshops and open-problem discussions.
\item Supplied literature source labeled ``Rogalski--Sierra'': \url{https://arxiv.org/abs/1101.0737}. The record supplies no further bibliographic details or specific supporting passage.
\item Literature assessment dated exactly \texttt{2026-08-21}: the supplied curation labels the problem open, describes both assertions as conjectures in the source, and reports that no general primary resolution was found. Its listed checked sources are the two links above. This is supplied dated curation, not independently verified literature assessment; the background explicitly requests independent verification of time-sensitive status and cited claims.
\item The registry labels the statement \texttt{corrected\_verified} and reports that it was checked against the official Oberwolfach report. These are supplied verification claims; this draft does not claim to have visited or independently checked either external source.
\end{itemize}
\end{document}
